% GENERATED FILE. Edit canonical lessons, then run npm run generate:books.
% canonical-source-hash: fnv1a64:1acc45fb5297827c
% canonical-lessons: ZH-C425-zhusu, ZH-C425-jiudian, ZH-W425-jiu-wine, ZH-C425-kefang, ZH-C425-fangka, ZH-C425-datang

\chapter{The Hotel}
\label{ch:zh-the-hotel}
\hlchaptermodality{425}

\begin{chapteropening}
\textbf{By the end of this chapter:} \emph{I can talk about a place to stay, a big hotel, its guest rooms, the room card and the lobby.}

Complete the last lesson of chapter 425: I can talk about a place to stay, a big hotel, its guest rooms, the room card and the lobby.
\end{chapteropening}

\section[zhùsù]{zhùsù — accommodation; a place to stay}

\label{lesson:ZH-C425-zhusu}



\begin{quote}
Before the new one: say the Mandarin for speed per hour, then the Mandarin for halfway.
\end{quote}

\subsection*{You'll want to know: zhùsù}
\emph{zhùsù} — \textquotedblleft{}accommodation; a place to stay\textquotedblright{}.

\begin{quote}
\emph{zhèr de zhùsù hěn guì} — Accommodation is expensive here.
\end{quote}

Its characters come later; for now it is a word for the ear.

\begin{cousinweb}
\emph{zhùsù} is built from \emph{zhù} 'to stay' and \emph{sù} 'to lodge for the night', both inherited from Middle Chinese, \emph{drjuH} and \emph{sjuwk} (Guangyun, in Baxter's transcription), which glosses the first 'to stop' and the second 'a lodging'.
\end{cousinweb}

\begin{grammarlens}[title={What you've built}]
One more word for getting around.
\end{grammarlens}

\subsection*{Your turn}
\begin{itemize}
\raggedright
  \item \emph{Say it:} \emph{zhùsù}
  \item \emph{Say it:} \emph{zhùsù}, once more
  \item \emph{Say it:} \emph{zhèr de zhùsù hěn guì} — Accommodation is expensive here.
  \item \emph{Recall:} say the Mandarin for speed per hour, then the Mandarin for halfway, then say \emph{zhùsù} again
\end{itemize}

\begin{tcolorbox}[breakable,colback=teal!4,colframe=teal!35!black,title={Before you move on}]
What does zhùsù mean? (\textquotedblleft{}Accommodation; a place to stay\textquotedblright{}.) Say it once more.
\end{tcolorbox}
\section[\texorpdfstring{\zh{酒店} (jiǔdiàn)}{酒店 (jiǔdiàn)}]{\zh{酒店} (jiǔdiàn) — big hotel}

\label{lesson:ZH-C425-jiudian}



\begin{quote}
Before the new one: say the Mandarin for halfway, then the Mandarin for a place to stay.
\end{quote}

\subsection*{You'll want to know: \zh{酒店}}
\textbf{\zh{酒店}} — \emph{jiǔdiàn} — \textquotedblleft{}big hotel\textquotedblright{}.

\begin{quote}
\textbf{\zh{这家酒店很大}} \emph{zhè jiā jiǔdiàn hěn dà} — This hotel is very big.
\end{quote}

The new character, in pieces: \zh{酒} — ten strokes in two pieces: on the left the water piece of \zh{洗}, three strokes; on the right a bar on top, a frame, a short stroke falling to the left and a vertical bending right inside it, a short bar, and a bar closing the foot. The next lesson puts it in your hand.

\begin{cousinweb}
\textbf{\zh{酒店}} (\emph{jiǔdiàn}) is built from \zh{酒} 'wine' and \zh{店} 'a shop, an inn', both inherited from Middle Chinese, \emph{tsjuwX} and \emph{temH} (Guangyun, in Baxter's transcription), which glosses the first as wine and sweet wine; how the pair came to mean a big hotel is not settled.
\end{cousinweb}

\begin{grammarlens}[title={What you've built}]
One more word for getting around.
\end{grammarlens}

\subsection*{Your turn}
\begin{itemize}
\raggedright
  \item \emph{Say it:} \emph{jiǔdiàn}
  \item \emph{Say it:} \emph{jiǔdiàn}, once more
  \item \emph{Say it:} \emph{zhè jiā jiǔdiàn hěn dà} — This hotel is very big.
  \item \emph{Recall:} say the Mandarin for halfway, then the Mandarin for a place to stay, then say \emph{jiǔdiàn} again
\end{itemize}

\begin{tcolorbox}[breakable,colback=teal!4,colframe=teal!35!black,title={Before you move on}]
What does \zh{酒店} mean? (\textquotedblleft{}Big hotel\textquotedblright{}.) Say it once more.
\end{tcolorbox}
\section[\texorpdfstring{\zh{酒} (jiǔ)}{酒 (jiǔ)}]{\zh{酒} — the water piece, then a jar}

\label{lesson:ZH-W425-jiu-wine}



\begin{quote}
The word first, then the character inside it.

\begin{itemize}
\raggedright
  \item \emph{Recall:} say \emph{jiǔdiàn} — \textbf{R1}, one lesson back, before it decays
\end{itemize}

\begin{quote}
\zh{洗} \zh{店}
\end{quote}

\emph{Read it:} each one aloud, with nothing said first — \emph{xǐ}, \emph{diàn}
\end{quote}

\subsection*{Script — the water piece, then a jar}
\hlblockfigure{\detokenize{figures/ZH-W425-jiu-wine-filmstrip.pdf}}{How \zh{酒} is written, stroke by stroke}

\begin{quote}
\zh{酒}
\end{quote}

\textbf{Two pieces, left before right.} Count the pieces, not the strokes.

\textbf{The left piece, the water piece, three strokes}, as in \zh{洗}: two dots, then a stroke that rises to the right.

\textbf{The right piece, seven strokes.} A bar across the top. A vertical down the left, and a bar that turns down the right side. Inside, a short stroke falling to the left and a vertical that bends to the right at the foot. A short bar across them. Then a bar that closes the foot.

3 + 7 = ten.

\subsection*{Your turn}
\begin{itemize}
\raggedright
  \item \emph{Write it:} \zh{酒}, all ten, with the model in front of you
  \item \emph{Cover:} write it again from memory, and count the pieces, not the strokes
\end{itemize}

\begin{tcolorbox}[breakable,colback=teal!4,colframe=teal!35!black,title={Before you move on}]
How many strokes in \zh{酒}? (\textbf{Ten}.) What is the left piece? (\textbf{The water piece, as in \zh{洗}}.) Is the frame closed at the foot? (\textbf{Yes, the last bar closes it}.)

Source: \href{https://raw.githubusercontent.com/chanind/hanzi-writer-data/68d10a4b21150cae5e1ebbd223eed289cf32d90c/data/\%E9\%85\%92.json}{Hanzi Writer Data \zh{酒}, pinned stroke paths and medians}.
\end{tcolorbox}
\section[\texorpdfstring{\zh{客房} (kèfáng)}{客房 (kèfáng)}]{\zh{客房} (kèfáng) — hotel guest room}

\label{lesson:ZH-C425-kefang}



\begin{quote}
Before the new one: say the Mandarin for a place to stay, then the Mandarin for a big hotel.
\end{quote}

\subsection*{You'll want to know: \zh{客房}}
\textbf{\zh{客房}} — \emph{kèfáng} — \textquotedblleft{}hotel guest room\textquotedblright{}.

\begin{quote}
\textbf{\zh{这家酒店有一百间客房}} \emph{zhè jiā jiǔdiàn yǒu yìbǎi jiān kèfáng} — This hotel has a hundred guest rooms.
\end{quote}

\begin{cousinweb}
\textbf{\zh{客房}} (\emph{kèfáng}) is built from \zh{客} 'a guest' and \zh{房} 'a room', both inherited from Middle Chinese, \emph{khaek} and \emph{bjang} (Guangyun, in Baxter's transcription).
\end{cousinweb}

\begin{grammarlens}[title={What you've built}]
One more word for getting around.
\end{grammarlens}

\subsection*{Your turn}
\begin{itemize}
\raggedright
  \item \emph{Say it:} \emph{kèfáng}
  \item \emph{Say it:} \emph{kèfáng}, once more
  \item \emph{Say it:} \emph{zhè jiā jiǔdiàn yǒu yìbǎi jiān kèfáng} — This hotel has a hundred guest rooms.
  \item \emph{Recall:} say the Mandarin for a place to stay, then the Mandarin for a big hotel, then say \emph{kèfáng} again
\end{itemize}

\begin{tcolorbox}[breakable,colback=teal!4,colframe=teal!35!black,title={Before you move on}]
What does \zh{客房} mean? (\textquotedblleft{}Hotel guest room\textquotedblright{}.) Say it once more.
\end{tcolorbox}
\section[fángkǎ]{fángkǎ — room key card}

\label{lesson:ZH-C425-fangka}



\begin{quote}
Before the new one: say the Mandarin for a big hotel, then the Mandarin for a hotel guest room.
\end{quote}

\subsection*{You'll want to know: fángkǎ}
\emph{fángkǎ} — \textquotedblleft{}room key card\textquotedblright{}.

\begin{quote}
\emph{wǒ de fángkǎ bú jiàn le} — My room card has gone missing.
\end{quote}

Its characters come later; for now it is a word for the ear.

\begin{cousinweb}
\emph{fángkǎ} is built from \zh{房} 'a room' and \emph{kǎ} 'a card'. The first is inherited from Middle Chinese \emph{bjang} (Guangyun, in Baxter's transcription); \emph{kǎ} is usually traced to English \emph{card}.
\end{cousinweb}

\begin{grammarlens}[title={What you've built}]
One more word for getting around.
\end{grammarlens}

\subsection*{Your turn}
\begin{itemize}
\raggedright
  \item \emph{Say it:} \emph{fángkǎ}
  \item \emph{Say it:} \emph{fángkǎ}, once more
  \item \emph{Say it:} \emph{wǒ de fángkǎ bú jiàn le} — My room card has gone missing.
  \item \emph{Recall:} say the Mandarin for a big hotel, then the Mandarin for a hotel guest room, then say \emph{fángkǎ} again
\end{itemize}

\begin{tcolorbox}[breakable,colback=teal!4,colframe=teal!35!black,title={Before you move on}]
What does fángkǎ mean? (\textquotedblleft{}Room key card\textquotedblright{}.) Say it once more.
\end{tcolorbox}
\section[dàtáng]{dàtáng — hotel lobby}

\label{lesson:ZH-C425-datang}



\begin{quote}
Before the new one: say the Mandarin for a hotel guest room, then the Mandarin for a room key card.
\end{quote}

\subsection*{You'll want to know: dàtáng}
\emph{dàtáng} — \textquotedblleft{}hotel lobby\textquotedblright{}.

\begin{quote}
\emph{wǒ zài jiǔdiàn de dàtáng děng nǐ} — I'll wait for you in the hotel lobby.
\end{quote}

Its characters come later; for now it is a word for the ear.

\begin{cousinweb}
\emph{dàtáng} is built from \zh{大} 'big' and \emph{táng} 'a hall', both inherited from Middle Chinese, \emph{daH} and \emph{dang} (Guangyun, in Baxter's transcription).
\end{cousinweb}

\begin{grammarlens}[title={What you've built}]
One more word for getting around.
\end{grammarlens}

\subsection*{Your turn}
\begin{itemize}
\raggedright
  \item \emph{Say it:} \emph{dàtáng}
  \item \emph{Say it:} \emph{dàtáng}, once more
  \item \emph{Say it:} \emph{wǒ zài jiǔdiàn de dàtáng děng nǐ} — I'll wait for you in the hotel lobby.
  \item \emph{Recall:} say the Mandarin for a hotel guest room, then the Mandarin for a room key card, then say \emph{dàtáng} again
\end{itemize}

\begin{tcolorbox}[breakable,colback=teal!4,colframe=teal!35!black,title={Before you move on}]
What does dàtáng mean? (\textquotedblleft{}Hotel lobby\textquotedblright{}.) Say it once more.
\end{tcolorbox}
